\RequirePackage{plautopatch}
\documentclass[platex,dvipdfmx,a4paper]{jsarticle}

\usepackage{amsmath}
\usepackage{amsfonts}
\usepackage{amssymb}

\usepackage[dvipdfmx]{graphicx}
\usepackage{float}

\usepackage{tikz}
\usepackage{caption}

% 「%」は以降の内容を「改行コードも含めて」無視するコマンド
\usepackage[%
  dvipdfmx,% 欧文ではコメントアウトする
  setpagesize=false,%
  bookmarks=true,%
  bookmarksdepth=tocdepth,%
  bookmarksnumbered=true,%
  colorlinks=false,%
  pdftitle={},%
  pdfsubject={},%
  pdfauthor={},%
  pdfkeywords={}%
]{hyperref}
% PDFのしおり機能の日本語文字化けを防ぐ((u)pLaTeXのときのみかく)
\usepackage{pxjahyper}

\usepackage{framed}
\setlength{\OuterFrameSep}{0pt}

\newcommand{\hb}[2]{\frac{\partial #1}{\partial #2}}
\newcommand{\bbn}[2]{\frac{d #1}{d #2}}
\newcommand{\nhb}[3]{\frac{\partial^{#1} #2}{\partial #3^{#1}}}
\newcommand{\nbbn}[3]{\frac{d^{#1} #2}{d #3^{#1}}}
\newcommand{\vectxy}[2]{\begin{bmatrix}#1\\#2\end{bmatrix}}
\newcommand{\vectxyz}[3]{\begin{bmatrix}#1\\#2\\#3\end{bmatrix}}
\newcommand{\vect}[1]{\boldsymbol{\mathrm{#1}}}

\newcommand{\ssum}[3]{\sum_{#1=#2}^#3}

\newcommand{\veff}{V_{\mathrm{eff}}(r)}

\hypersetup{colorlinks=true}
\hypersetup{linkcolor=red}

\title{解析力学演習レポート2}
\author{25B00729　嶋田 悠音}

\begin{document}
\maketitle
\subsection*{1}

スカラー場 $\phi(\vect r,t)$ が

\begin{align*}
\left(
\nabla^2-\frac{1}{c^2}\frac{\partial^2}{\partial t^2}
\right)\phi(\vect r,t)
=
\rho(\vect r,t)
\end{align*}

を満たすとする。

まずグリーン関数 $G(\vect r,t;\vect r',t')$ を

\begin{align*}
\left(
\nabla^2-\frac{1}{c^2}\frac{\partial^2}{\partial t^2}
\right)
G(\vect r,t;\vect r',t')
=
\delta^{(3)}(\vect r-\vect r')
\delta(t-t')
\end{align*}

によって定義する。

空間および時間の並進対称性より、

\begin{align*}
G(\vect r,t;\vect r',t')
=
G(\vect r-\vect r',t-t')
\end{align*}

である。

このとき波動方程式の特解は

\begin{align*}
\phi(\vect r,t)
=
\int d^3r'
\int dt'\,
G(\vect r-\vect r',t-t')
\rho(\vect r',t')
\end{align*}

と表される。

\subsection*{2}

グリーン関数のフーリエ変換

\begin{align*}
G(\vect r-\vect r',t-t')
=
\frac{1}{(2\pi)^4}
\int d^3k
\int d\omega\,
\widetilde G(\vect k,\omega)
e^{i\vect k\cdot(\vect r-\vect r')-i\omega(t-t')}
\end{align*}

を導入する。

またデルタ関数のフーリエ表示

\begin{align*}
\delta^{(3)}(\vect r-\vect r')
\delta(t-t')
=
\frac{1}{(2\pi)^4}
\int d^3k
\int d\omega\,
e^{i\vect k\cdot(\vect r-\vect r')-i\omega(t-t')}
\end{align*}

を用いる。

これらをグリーン関数の定義式へ代入すると、

\begin{align*}
\left(
-k^2+\frac{\omega^2}{c^2}
\right)
\widetilde G(\vect k,\omega)
=
1
\end{align*}

を得る。

したがって

\begin{align*}
\widetilde G(\vect k,\omega)
=
\frac{1}
{\omega^2/c^2-k^2}
=
-\frac{1}
{k^2-\omega^2/c^2}
\end{align*}

である。

よって

\begin{align*}
G(\vect r-\vect r',t-t')
=
-\frac{1}{(2\pi)^4}
\int d^3k
\int d\omega\,
\frac{
e^{i\vect k\cdot(\vect r-\vect r')-i\omega(t-t')}
}{
k^2-\omega^2/c^2
}
\end{align*}

となる。

\subsection*{3}

遅延ポテンシャルを求めるため、

\begin{align*}
G(\vect r-\vect r',t-t')
=
0
\qquad
(t<t')
\end{align*}

を課す。

この条件を実現するため、

\begin{align*}
G_{\mathrm{ret}}(\vect r-\vect r',t-t')
=
-\frac{1}{(2\pi)^4}
\int d^3k
\int d\omega\,
\frac{
e^{i\vect k\cdot(\vect r-\vect r')-i\omega(t-t')}
}{
k^2-(\omega+i0)^2/c^2
}
\end{align*}

を考える。

分母を因数分解すると、

\begin{align*}
k^2-\frac{(\omega+i0)^2}{c^2}
=
-\frac{1}{c^2}
(\omega-ck+i0)
(\omega+ck+i0)
\end{align*}

であるから、

\begin{align*}
G_{\mathrm{ret}}
=
\frac{c^2}{(2\pi)^4}
\int d^3k
\int d\omega\,
\frac{
e^{i\vect k\cdot(\vect r-\vect r')-i\omega(t-t')}
}{
(\omega-ck+i0)
(\omega+ck+i0)
}
\end{align*}

と書ける。

$t>t'$ とする。このとき指数因子

\begin{align*}
e^{-i\omega(t-t')}
\end{align*}

のため、積分路を下半平面で閉じると円弧部分の寄与は消える。

したがって留数定理より、

\begin{align*}
\int_{-\infty}^{\infty}
d\omega\,
\frac{
e^{-i\omega(t-t')}
}{
(\omega-ck+i0)
(\omega+ck+i0)
}
=
-\frac{\pi i}{ck}
\left[
e^{-ick(t-t')}
-
e^{ick(t-t')}
\right]
\end{align*}

を得る。

これを代入すると、

\begin{align*}
G_{\mathrm{ret}}
=
\Theta(t-t')
\frac{c}{(2\pi)^3}
\int d^3k\,
e^{i\vect k\cdot(\vect r-\vect r')}
\frac{\sin\!\left[ck(t-t')\right]}{k}
\end{align*}

となる。

\subsection*{4}

ここで

\begin{align*}
R=|\vect r-\vect r'|
\end{align*}

とおく。

$\vect r-\vect r'$ の方向を極軸に選ぶと、

\begin{align*}
\int d\Omega\,
e^{i\vect k\cdot(\vect r-\vect r')}
=
4\pi
\frac{\sin(kR)}{kR}
\end{align*}

であるから、

\begin{align*}
G_{\mathrm{ret}}
=
\Theta(t-t')
\frac{c}{2\pi^2R}
\int_0^\infty
dk\,
\sin(kR)
\sin\!\left[ck(t-t')\right]
\end{align*}

を得る。

積和公式

\begin{align*}
\sin A\sin B
=
\frac12
\left[
\cos(A-B)-\cos(A+B)
\right]
\end{align*}

を用いると、

\begin{align*}
G_{\mathrm{ret}}
=
\Theta(t-t')
\frac{c}{4\pi^2R}
\int_0^\infty
dk
\left[
\cos\!\left(k(R-c(t-t'))\right)
-
\cos\!\left(k(R+c(t-t'))\right)
\right]
\end{align*}

となる。

さらに

\begin{align*}
\int_0^\infty \cos(ak)\,dk
=
\pi\delta(a)
\end{align*}

を用いると、

\begin{align*}
G_{\mathrm{ret}}
=
\Theta(t-t')
\frac{c}{4\pi R}
\left[
\delta\!\left(R-c(t-t')\right)
-
\delta\!\left(R+c(t-t')\right)
\right]
\end{align*}

を得る。

$t>t'$ および $R>0$ のもとでは

\begin{align*}
\delta\!\left(R+c(t-t')\right)=0
\end{align*}

なので、

\begin{align*}
G_{\mathrm{ret}}
=
\Theta(t-t')
\frac{c}{4\pi R}
\delta\!\left(R-c(t-t')\right)
\end{align*}

となる。

最後に

\begin{align*}
\delta\!\left(R-c(t-t')\right)
=
\frac1c
\delta\!\left(
t-t'-\frac{R}{c}
\right)
\end{align*}

を用いると、

\begin{align*}
G_{\mathrm{ret}}(\vect r-\vect r',t-t')
=
\frac{1}{4\pi|\vect r-\vect r'|}
\delta\!\left(
t-t'
-\frac{|\vect r-\vect r'|}{c}
\right)
\end{align*}

を得る。

\subsection*{5}

これを特解の表式へ代入する：

\begin{align*}
\phi(\vect r,t)
&=
\int d^3r'
\int dt'\,
\frac{
\delta\!\left(
t-t'
-\frac{|\vect r-\vect r'|}{c}
\right)
}
{4\pi|\vect r-\vect r'|}
\rho(\vect r',t') \\
&=
\frac{1}{4\pi}
\int d^3r'
\frac{
\rho\!\left(
\vect r',
t-\frac{|\vect r-\vect r'|}{c}
\right)
}
{|\vect r-\vect r'|}。
\end{align*}

したがって、波動方程式

\begin{align*}
\left(
\nabla^2-\frac{1}{c^2}\frac{\partial^2}{\partial t^2}
\right)\phi
=
\rho
\end{align*}

の遅延ポテンシャルは

\begin{align*}
\phi(\vect r,t)
=
\frac{1}{4\pi}
\int d^3r'
\frac{
\rho\!\left(
\vect r',
t-\frac{|\vect r-\vect r'|}{c}
\right)
}
{|\vect r-\vect r'|}
\end{align*}

で与えられる。

\subsection*{6}

一般解は、上で求めた特解に斉次方程式

\begin{align*}
\left(
\nabla^2-\frac{1}{c^2}\frac{\partial^2}{\partial t^2}
\right)\phi_{\mathrm{hom}}
=
0
\end{align*}

の解を加えたものである。

しかし、遅延条件

\begin{align*}
G(\vect r-\vect r',t-t')
=
0
\qquad
(t<t')
\end{align*}

および無限遠からの入射波が存在しない放射境界条件を課すと、斉次解は存在しない。

したがって遅延ポテンシャルは上式によって一意に定まり、その値は過去のソース分布のみで決定される。

以上より、グリーン関数法とフーリエ変換を用いて波動方程式の遅延ポテンシャルが導出された。
\end{document}